\section{Perturbation Operator: Formal Definition}\label{sec:operator}

We now define formally the object on which the kernel implementation
acts.

\subsection{State Space and Perturbation Algebra}

We adopt the BEAMFS-compatible state space
\(\mathcal{S} = \mathcal{S}_V \times \mathcal{S}_P\)
of~\cite{desbrieres2026beamfsv1}, where \(\mathcal{S}_V\) is the
volatile component (kernel data structures, page cache, in-memory
mirrors) and \(\mathcal{S}_P\) is the persistent component (on-disk
data and metadata).

\begin{definition}[Perturbation operator]\label{def:perturbation}
The RadFI perturbation operator \(\varepsilon\) is a partial function
\(\varepsilon : \mathcal{S} \times \Pi \to \mathcal{S}\),
parameterized by a perturbation parameter set
\(\Pi = \langle \mathit{filters}, \mathit{prob}, \mathit{seed} \rangle\),
and defined by the following procedure on input \((s, \pi)\):
\begin{enumerate}
\item Intercept a bio submission \(b\) in transit through
  \texttt{submit\_bio\_noacct} or \texttt{submit\_bh}.
\item Apply the filter algebra \(\mathit{filters}(\pi)\) to \(b\).
  If \(b\) is rejected, return \(s\) unchanged.
\item Sample the PRNG seeded by \(\mathit{seed}(\pi)\). If the
  sample exceeds \(\mathit{prob}(\pi)\), return \(s\) unchanged.
\item Otherwise, select a byte in the first non-empty
  \texttt{bio\_vec} segment of \(b\) uniformly at random, and flip
  one bit at a uniformly sampled position \(0 \leq p < 8\).
\item Return the modified state \(s'\).
\end{enumerate}
\end{definition}

\begin{definition}[Filter algebra]\label{def:filter}
The filter set \(\mathit{filters}\) is a tuple
\(\langle \mathit{dev}, \mathit{inode}, \mathit{block}, \mathit{op} \rangle\)
where each component is either a wildcard (admits any value) or a
concrete value. A bio \(b\) passes the filter iff every non-wildcard
component matches the corresponding attribute of \(b\):
\begin{itemize}
\item \(\mathit{dev}\): the encoded device number
  \texttt{bio->bi\_bdev->bd\_dev} (kernel encoding
  \texttt{(major << 8) | minor} for legacy minors).
\item \(\mathit{inode}\): the inode number, available only on the
  \texttt{submit\_bh} hook path where the buffer head carries
  \texttt{b\_assoc\_map}.
\item \(\mathit{block}\): the sector number
  \texttt{bh->b\_blocknr}, available on the \texttt{submit\_bh} path.
\item \(\mathit{op}\): the bio operation type, restricted to
  \(\{\texttt{REQ\_OP\_READ}, \texttt{REQ\_OP\_WRITE}\}\) in v0.1.2;
  other ops are filtered out.
\end{itemize}
\end{definition}

The inode and block filters are bound to the \texttt{submit\_bh}
hook because a \texttt{bio} may aggregate buffers from multiple
inodes; RadFI v0.1.x does not attempt fine-grained per-buffer
targeting on the bio path.

\subsection{Probability Sampling and PRNG}

The probability \(\mathit{prob}\) is expressed in parts per million
(PPM): \(\mathit{prob} = 10^6\) flips every intercepted bio that
passes the filter; \(\mathit{prob} = 10^3\) flips one bio in a
thousand. The PRNG is a xoshiro256** generator~\cite{blackman2021xoshiro},
seeded by the user-supplied \(\mathit{seed}\). A fixed seed produces
a deterministic flip sequence, enabling regression on resilience
properties: a recovery theorem that holds for seed \(S_1\) but not
for \(S_2\) is a falsification, regardless of the absolute flip
count.

\subsection{The Five Invariants of a Faithful SEU Emulator}\label{ssec:invariants}

The faithfulness theorem (\cref{thm:faithfulness}) requires the
following invariants of the perturbation operator. They are minimal
and stated separately so that an implementation can be audited
against them line by line.

\begin{description}
\item[F1 -- Single-bit primitive.] Each invocation of
  \(\varepsilon\) that produces a perturbation flips exactly one
  bit. Multi-flip is by repeated invocation, not by single-call
  amplification. This matches the canonical single-event abstraction
  of~\cite{dodd2003seu,baumann2005softerrors}.
\item[F2 -- Targeting independence from data content.] The byte and
  bit position selected by \(\varepsilon\) are sampled from the bio
  payload's geometry, not from its contents. A flip at offset 17 in
  a bio is not statistically correlated with the value at offset 17.
\item[F3 -- Filter monotonicity.] Adding a non-wildcard filter
  component to \(\mathit{filters}\) can only reduce the set of
  intercepted bios, never enlarge it. A filtered injection is always
  a sub-distribution of an unfiltered one.
\item[F4 -- Replay determinism.] Given fixed
  \((\mathit{filters}, \mathit{prob}, \mathit{seed})\) and a fixed
  I/O sequence, \(\varepsilon\) produces the same flip sequence
  across runs.
\item[F5 -- Counter accounting.] \(\varepsilon\) maintains atomic
  counters \((\mathit{call\_count}, \mathit{flip\_count},
  \mathit{skipped\_disabled}, \mathit{skipped\_prob})\) such that
  the sum
  \(\mathit{flip\_count} + \mathit{skipped\_disabled} + \mathit{skipped\_prob}\)
  equals the number of post-filter intercepted bios.
\end{description}

F1 grounds the SEU semantics. F2 is the absence-of-bias property.
F3 is the algebraic well-formedness of the filter lattice. F4 is
what makes RadFI a falsification instrument rather than a stochastic
noise source. F5 is the audit surface that allows an empirical run
to be self-validated against the expected statistics.
